\documentclass[a4paper,11pt]{letter}

\usepackage[left=15mm,right=15mm,top=10mm,bottom=5mm]{geometry}

\usepackage{hyperref}

\usepackage{fontspec}
\setmainfont{Fira Sans}
\linespread{1.05}

\usepackage{xcolor} % For custom colors
\usepackage{fontawesome5} % For email and phone symbols

\hypersetup{
    colorlinks=true,
    linkcolor=blue,
    urlcolor=blue
}

\newcommand{\senderinfo}{
    \noindent
    \begin{minipage}{0.45\textwidth}
    \raggedleft
    \textbf{Sophie Lund Wagner} \\
    Falen 12 4.th., 5000 Odense, Denmark \\
    \faEnvelope\ \href{mailto:sowag21@student.sdu.dk}{sowag21@student.sdu.dk} \\
    \faPhone\ \href{tel:+45 42187737}{+45 42187737} \\
    \end{minipage}
}

 % Sender's name for closing
\date{\today}


\begin{document}

\pagestyle{empty}

\begin{flushright}
\textbf{Sophie Lund Wagner } \\
   Falen 12 4.th., 5000 Odense, Denmark \\
    \faEnvelope\ \href{mailto:sowag21@student.sdu.dk}{sowag21@student.sdu.dk} \\
    \faPhone\ \href{tel:+45 42187737}{+45 42187737} \\
\end{flushright}

Cosmology, Particle Astrophysics and Strings division, \\
Oskar Klein Centre, \\
Stockholm, Sweden \\

Dear Members of the Selection Committee,

My interest in theoretical physics has been driven by a desire to understand the fundamental structure of nature, which has naturally led me towards quantum field theory, cosmology and particle physics. \textbf{I am particularly fascinated by the role of symmetries in shaping physical laws}, including how phenomena such as spontaneous symmetry breaking and anomalies play a central role in determining particle behaviour, and how these mechanisms may leave observable imprints in cosmology through dark matter or gravitational waves. I wish to deepen my understanding of these principles and ultimately piece the clues together to answer the question: \textbf{what symmetries govern physics beyond the Standard Model, and how did they drive the evolution of the early universe?} 

\textbf{My research experience spans fundamental particle physics, cosmology, and quantum field theory. Currently, I am working on my master’s thesis}, which focuses on gaining a deeper understanding of fundamental interactions and \textbf{beyond-the-Standard-Model physics.} In my research, I explore grand unified theories (GUTs). Rather than proposing a specific gauge group, \textbf{I study the landscape of GUTs, and aim to extend this approach to include scalars}, by exploring how their representations affect model viability. This work is currently being prepared for publication. In parallel, I am involved in a project on infinite heat order, concerning \textbf{high-temperature symmetry non-restoration in UV-complete models.} During my master’s degree, while on exchange in Canada, I have completed another project on topological anomalies in quantum field theory. \\
During my bachelor's degree, I worked on density profiles of dark matter, both theoretically and using machine learning techniques. \textbf{This interest culminated in a bachelor’s thesis in cosmology,} where I examined whether inhomogeneities in the structure of the universe could explain its accelerated expansion.

For my PhD, I aim to study the interplay between quantum field theory and cosmology. I aim to apply the framework of quantum field theory to study symmetry breaking, phase transitions, and beyond-Standard-Model physics, with the hope of uncovering the mechanisms governing fundamental interactions. I am
\textbf{motivated by open questions at the interface of particle physics and cosmology, such as the nature of dark matter, the origin of cosmic structure, and how new physics can be probed through cosmological observations.} Furtermore, the possibility of glimpsing information about the early universe through effective field theory methods excites me. Ultimately, I aim
to help clarify the underlying principles that determine the structure of our universe, \textbf{in close alignment with the research directions pursued by the \textit{Cosmology, Particle Astrophysics and Strings division}.} 

I am excited to apply for the PhD position in the Cosmology, Particle Astrophysics and Strings division at Stockholm University. The group’s work on early-universe cosmology, phase transitions, and particle physics beyond the Standard Model strongly aligns with my background. \textbf{I am particularly drawn to the interdisciplinary environment of the Oskar Klein Centre, where theoretical developments are directly informed by current experiments and astrophysical observations.} The opportunity to contribute to this collaborative environment at the interface of high-energy phenomenology and particle cosmology would be deeply motivating. 

Thank you for considering my application. I would be thrilled to contribute to the research activities of the Cosmology, Particle Astrophysics and Strings division.

Sincerely,

Sophie Lund Wagner 
\newpage
\vspace{5cm}
\begin{center}
    {\Large \textbf{List of Referees}} \\[1em]
\end{center}

\vspace{2em}

\noindent
\textbf{Referee 1} \\[0.3em]
\textbf{Francesco Sannino} \\
Director, Quantum Theory Center ($\hbar$QTC) and Chair of Science\\
Danish Institute for Advanced Study (Danish IAS) \\
University of Southern Denmark \\
Campusvej 55, DK-5230 Odense M, Denmark \\[0.4em]
Professor of Theoretical Physics\\
Department of Physics E.\ Pancini, Università di Napoli Federico II \\
INFN Sezione di Napoli \\
Scuola Superiore Meridionale \\
Via Cintia, 80126 Napoli, Italy\\[0.4em]
\faEnvelope\ \href{mailto:sannino@qtc.sdu.dk}{sannino@qtc.sdu.dk} \\
\faEnvelope\ \href{mailto:sannino@mac.com}{sannino@mac.com} \\

\vspace{2em}

\noindent
\textbf{Referee 2} \\[0.3em]
\textbf{ Giacomo Cacciapaglia} \\
Senior Researcher \\
Laboratoire de Physique Théorique et Hautes Énergies (LPTHE), UMR 7589 \\
Sorbonne Université \& CNRS \\
4 place Jussieu, 75252 Paris Cedex 05, France\\[0.4em]
\faEnvelope\ \href{cacciapa@lpthe.jussieu.fr}{cacciapa@lpthe.jussieu.fr} 

\vspace{2em}

% Uncomment if you want a third referee
\noindent
\textbf{Referee 3} \\[0.3em]
\textbf{Sofie Marie Koksbang} \\
Associate Professor \\
CP3-Origins \\
University of Southern Denmark \\
Campusvej 55, DK-5230 Odense M, Denmark \\[0.4em]
\faEnvelope\ \href{mailto:koksbang@cp3.sdu.dk}{koksbang@cp3.sdu.dk} 
\end{document}  